\documentclass[10pt]{article}
\usepackage[T1]{fontenc}
\usepackage[utf8]{inputenc}
\usepackage[french]{babel}
\usepackage{amsmath,amssymb,mathrsfs,mathpazo}
\usepackage{geometry}
\geometry{a4paper,margin=2cm}
\usepackage{graphicx,array,tabularx,tabularray,colortbl,multirow,multicol,booktabs,enumitem,pifont,ulem,fancybox,diagbox,cancel}
\usepackage{pstricks,pstricks-add,pst-plot,pst-tree,pst-node,pst-text,pst-func,pst-math,pst-bspline,pst-3dplot,pst-eucl,pst-solides3d}
\let\pscalculate\undefined
\usepackage{pst-bezier}
\usepackage{tikz,pgfplots,tkz-tab}
\usetikzlibrary{arrows,arrows.meta,patterns,calc,positioning,trees,decorations.markings,shapes.geometric,shapes.misc,shapes.arrows,intersections,decorations.pathmorphing,decorations.pathreplacing,backgrounds,shadows,babel}
\pgfplotsset{compat=1.18}
\usepgfplotslibrary{fillbetween}
\usepackage[np]{numprint}
\usepackage{listings,makecell,pas-tableur,twoopt}
\let\barre\undefined
\usepackage[pstricks]{tablvar}
\definecolor{darkgreen}{RGB}{0,100,0}
\providecommand{\red}{\color{red}}
\providecommand{\blue}{\color{blue}}
\providecommand{\green}{\color{green}}
\providecommand{\black}{\color{black}}
\providecommand{\white}{\color{white}}
\providecommand{\gray}{\color{gray}}
\providecommand{\cyan}{\color{cyan}}
\providecommand{\magenta}{\color{magenta}}
\providecommand{\yellow}{\color{yellow}}
\providecommand{\violet}{\color{violet}}
\providecommand{\orange}{\color{orange}}
\providecommand{\euro}{\texteuro}
\providecommand{\EUR}{\texteuro}
\providecommand{\ieme}{\textsuperscript{e}}
\providecommand{\degres}{\ensuremath{^\circ}}
\providecommand{\degree}{\ensuremath{^\circ}}
\providecommand{\vv}[1]{\overrightarrow{#1}}
\providecommand{\Coord}[1]{#1}
\providecommand{\decosix}{}
\providecommand{\starredbullet}{\ensuremath{\bullet}}
\pagestyle{empty}
\setlength{\parindent}{0pt}
\usepackage{adjustbox}\DeclareUnicodeCharacter{2236}{:}\DeclareUnicodeCharacter{2190}{\ensuremath{\leftarrow}}
\providecommand{\R}{\mathbb{R}}
\providecommand{\N}{\mathbb{N}}
\providecommand{\Z}{\mathbb{Z}}
\providecommand{\Q}{\mathbb{Q}}
\providecommand{\D}{\mathbb{D}}
\providecommand{\C}{\mathbb{C}}
\providecommand{\e}{{\mathrm{e}}}
\providecommand{\vect}[1]{\overrightarrow{#1}}
\providecommand{\vectt}[1]{\overrightarrow{\mathrm{#1}}}
\providecommand{\vvt}[1]{\overrightarrow{\mathrm{#1}}}
\providecommand{\np}[2][]{{#2~#1}}
\providecommand{\vv}[1]{{\overrightarrow{#1}}}
\renewcommand{\d}{{\mathrm{d}}}
\providecommand{\degres}{{^\circ}}
\providecommand{\degre}{{^\circ}}
\providecommand{\blue}{\color{blue}}
\providecommand{\red}{\color{red}}
\providecommand{\enskip}{\quad}
\providecommand{\llbracket}{\lbrack\!\lbrack}
\providecommand{\rrbracket}{\rbrack\!\rbrack}
\providecommand{\euro}{{\text{€}}}
\providecommand{\no}{{\text{n°}}}
\providecommand{\interval}[2]{{[#1;#2]}}
\providecommand{\intervFO}[2]{{[#1;#2[}}
\providecommand{\dotfill}{\ldots}

\providecommand{\ds}{\displaystyle}
\providecommand{\textding}[1]{#1}
\providecommand{\Oij}{\(\mathrm{O};\vec{\imath},\vec{\jmath}\)}
\providecommand{\Oijk}{\(\mathrm{O};\vec{\imath},\vec{\jmath},\vec{k}\)}
\providecommand{\Ouv}{\(\mathrm{O};\vec{u},\vec{v}\)}
\providecommand{\intervFF}[2]{\ensuremath{[#1;#2]}}
\renewcommand{\euro}{\ifmmode\text{\texteuro}\else\texteuro\fi}
\DeclareUnicodeCharacter{00B0}{\ifmmode^{\circ}\else\textdegree\fi}
\DeclareUnicodeCharacter{20AC}{\ifmmode\text{\texteuro}\else\texteuro\fi}
% Énoncés : Éducation nationale. Transcription et corrigés : APMEP, Denis Vergès et contributeurs. Découpage et classement : CoopMaths, Sébastien Lozano, Éric Elter et contributeurs. Adaptation de la présentation : Maths Explorer.
% CC BY-SA : https://coopmaths.fr/www/about/
% https://coopmaths.fr/alea/static/e3c/2020/tex/e3c_2020_00_sujet40_1.tex
\begin{document}
\noindent\textit{APMEP / CoopMaths. CC BY-SA. Présentation adaptée par Maths Explorer.}\par\medskip
\medskip

Ce QCM comprend 5 questions indépendantes. Pour chacune d'elles, une seule des réponses proposées est exacte.

Indiquer pour chaque question sur la copie la lettre correspondant à la réponse choisie.

Aucune justification n'est demandée.

Chaque réponse correcte rapporte $1$ point. Une réponse incorrecte ou une absence de réponse n'apporte ni ne retire de point.

\medskip

\begin{enumerate}
\item Pour tout réel $x$, \,$\sin(7\pi - x)$ est égal à :

\begin{center}
\begin{tabularx}{\linewidth}{|*{4}{X|}}\hline
\textbf{a.~~}$\sin x$&\textbf{b.~~}$- \sin x$&\textbf{c.~~}$\cos x$&\textbf{d.~~}$- \cos x$\\ \hline
\end{tabularx}
\end{center}

\item  Dans laquelle des quatre situations proposées ci-dessous le produit scalaire $\vect{\text{AB}} \cdot \vect{\text{AC}}$ est- il égal à 6 ?

\begin{center}
\begin{tabularx}{\linewidth}{|*{4}{>{\small}X|}}\hline
\textbf{a.~~}ABC est un triangle tel que: 

AB $= 6$, AC $= 4$ et 

BC $= 8$.
&\textbf{b.~~}Dans un repère orthonormé du plan: A$(-3~;~ 5)$, B$(2~;~-2)$
et C(1~;~7).
&\textbf{c.~~}ABC est un triangle rectangle en B tel que :
 
AB $= 3$ et $\text{BC} = 2$
&\textbf{d.~~}ABC est un triangle tel que :

AB $= 6$, AC $= 4$ et 

$\widehat{\text{BAC}} = 30\degres$.\\ \hline
\end{tabularx}
\end{center}
 
\item On considère la fonction $f$ définie sur $\R$ par: $f(x) = \dfrac{3x + 4}{x^2 + 1}$.

$f$ est dérivable sur $\R$ et, pour tout réel $x$,\; $f'(x)$ est égal à :

\begin{center}
\begin{tabularx}{\linewidth}{|*{4}{X|}}\hline
\textbf{a.~~}$\dfrac{3}{2x}$&\textbf{b.~~}$\dfrac{9x^2+8x+3\rule{0pt}{10pt}}{\left(x^2+1\right)^2}$&\textbf{c.~~}$\dfrac{-3x^2- 8x+3\rule{0pt}{10pt}}{\left(x^2+1\right)^2}$&\textbf{d.~~}$9x^2 + 8x+ 3$\\ \hline
\end{tabularx}
\end{center}

\item  Le plan est rapporté à un repère orthonormé.

L'ensemble des points $M(x~;~y)$ tels que $x^2 + y^2 - 10x +6y +30 = 0$ est :

\begin{center}
\begin{tabularx}{\linewidth}{|*{4}{X|}}\hline
\textbf{a.~~}une droite&\textbf{b.~~}une parabole&\textbf{c.~~}un cercle&\textbf{d.~~}ni une droite, ni une parabole, ni un cercle.\\ \hline
\end{tabularx}
\end{center}

\item La somme $1 + 5 + 5^2 + 5^3 + \ldots + 5^{30}$ est égale à :

\begin{center}
\begin{tabularx}{\linewidth}{|*{4}{X|}}\hline
\textbf{a.~~}$\dfrac{1 - 5^{30}\rule{0pt}{10pt}}{4\rule[-3pt]{0pt}{0pt}}$
&\textbf{b.~~}$\dfrac{5^{30} - 1}{4}$
&\textbf{c.~~}$\dfrac{1 - 5^{31}}{4}$
&\textbf{d.~~}$\dfrac{5^{31} - 1}{4}$\\ \hline
\end{tabularx}
\end{center}
\end{enumerate}

\bigskip

\end{document}
